\thispagestyle{empty}
\begin{center}
    \vspace*{1.2cm}
    {\Large\bfseries\scshape\color{focpgreendark} Remerciements}
    \vspace{0.8cm}
\end{center}

Au terme de ce Projet de Fin d'Année marquant une étape déterminante de ma formation d'ingénieur en Ingénierie Informatique et Réseaux, je tiens à exprimer ma profonde gratitude à toutes les personnes qui ont contribué, par leurs conseils, leur expertise et leur soutien bienveillant, à la réussite de ce travail.

Mes remerciements les plus sincères et respectueux s'adressent en premier lieu à mon tuteur pédagogique, \textbf{Saber Sarra}, pour son encadrement attentif, sa disponibilité constante, ses conseils méthodologiques avisés et ses orientations scientifiques précieuses. Ses retours constructifs et son exigence intellectuelle ont constitué un guide inestimable tout au long de la conception de ce rapport et du développement des solutions logicielles.

J'exprime ma vive reconnaissance à l'organisme d'accueil, la \textbf{Fondation OCP}, pour m'avoir offert l'opportunité d'effectuer ce projet de fin d'année au sein d'une institution de référence, engagée au service du développement durable et de l'innovation sociale au Maroc.

J'adresse mes chaleureux remerciements à l'ensemble de l'équipe professionnelle de la Fondation OCP, en particulier aux équipes de l'\textbf{Axe Social et Solidaire} et de la \textbf{Direction de la Transformation Digitale (DTD)}. Je les remercie pour la qualité de leur accueil, la clarté de leurs directives métier, la pertinence de leurs éclairages sur les problématiques de gestion des bénéficiaires et des coopératives, ainsi que pour leur collaboration active lors des séances d'élicitation et de validation des besoins.

Je tiens également à témoigner ma gratitude envers le corps professoral et l'équipe pédagogique de l'\textbf{École Marocaine des Sciences de l'Ingénieur (EMSI)} pour la rigueur et l'excellence de la formation scientifique et technique dispensée tout au long de mon cursus d'ingénieur. Les compétences acquises en architecture logicielle, en modélisation UML, en cybersécurité et en génie logiciel ont constitué le socle indispensable à la réalisation de ce projet.

Enfin, je remercie très cordialement toutes les personnes qui, de près ou de loin, ont apporté leur concours, leurs suggestions techniques ou leurs encouragements fraternels tout au long de cette aventure académique et professionnelle.

\vfill
\newpage
